\begin{prop}
    El producto directo finito de grupos solubles (resp. nilpotentes) es soluble (resp. nilpotente).
\end{prop}

\begin{proof}
    Sea $G = H \times K$ el producto directo de dos grupos. Probaremos el resultado para dos factores; el caso general se sigue por inducción matemática sobre el número de factores.
    
    \itshape 1. Caso Soluble: \normalfont
    Supongamos que $H$ y $K$ son solubles. Existen enteros $n_1, n_2$ tales que $H^{(n_1)} = \{e_H\}$ y $K^{(n_2)} = \{e_K\}$. Sea $n = \max(n_1, n_2)$.
    
    Recordemos que las operaciones en el producto directo se realizan componente a componente. El conmutador de dos elementos $(h_1, k_1)$ y $(h_2, k_2)$ en $G$ es:
    \begin{align*}
        [(h_1, k_1), (h_2, k_2)] &= (h_1, k_1)(h_2, k_2)(h_1, k_1)^{-1}(h_2, k_2)^{-1} \\
        &= (h_1 h_2 h_1^{-1} h_2^{-1}, k_1 k_2 k_1^{-1} k_2^{-1}) \\
        &= ([h_1, h_2], [k_1, k_2]).
    \end{align*}
    Por inducción sobre $m$, se cumple que el término $m$-ésimo de la serie derivada es el producto directo de los términos de los factores:
    \begin{equation*}
        (H \times K)^{(m)} = H^{(m)} \times K^{(m)}.
    \end{equation*}
    Evaluando en $n = \max(n_1, n_2)$:
    \begin{equation*}
        G^{(n)} = H^{(n)} \times K^{(n)} = \{e_H\} \times \{e_K\} = \{(e_H, e_K)\} = \{e_G\}.
    \end{equation*}
    Por lo tanto, $G$ es soluble.

    \itshape 2. Caso Nilpotente: \normalfont
    Supongamos que $H$ y $K$ son nilpotentes. Existen $n_1, n_2$ tales que $H_{(n_1)} = \{e_H\}$ y $K_{(n_2)} = \{e_K\}$. Sea $n = \max(n_1, n_2)$.
    
    De manera análoga, el conmutador de un elemento de $G$ con uno de un término de la serie central se comporta componente a componente:
    \begin{equation*}
        [ (g_h, g_k), (x_h, x_k) ] = ([g_h, x_h], [g_k, x_k]).
    \end{equation*}
    Esto implica que para la serie central descendente se cumple:
    \begin{equation*}
        (H \times K)_{(m)} = H_{(m)} \times K_{(m)}.
    \end{equation*}
    Evaluando en $n$:
    \begin{equation*}
        G_{(n)} = H_{(n)} \times K_{(n)} = \{e_H\} \times \{e_K\} = \{e_G\}.
    \end{equation*}
    Por lo tanto, $G$ es nilpotente.
\end{proof}

\ObservacionBox{Ejemplos y Notas}{
    \begin{enumerate}
        \item El grupo simétrico $S_n$ no es soluble para $n \ge 5$ (como vimos con $A_n$). Dado que todo grupo nilpotente es soluble, $S_n$ tampoco es nilpotente para $n \ge 5$.
        
        \item Sea $G$ un $p$-grupo finito (es decir, $|G| = p^n$ con $n \ge 1$). Entonces $G$ es nilpotente.
        
        \item Si $G$ es un producto directo de sus subgrupos de Sylow, entonces $G$ es nilpotente. Esto se sigue de la proposición anterior, ya que los $p$-grupos son nilpotentes y el producto directo de nilpotentes es nilpotente.
    \end{enumerate}
}

\begin{proof}[Demostración de que todo $p$-grupo es nilpotente]
    Procedemos por inducción sobre el orden de $G$, es decir, sobre el exponente $n$ donde $|G|=p^n$.
    
    \itshape Base de inducción ($n=1$): \normalfont
    Si $|G|=p$, entonces $G \cong \mathbb{Z}_p$, que es abeliano y por tanto nilpotente.
    
    \itshape Hipótesis de inducción: \normalfont
    Supongamos que todo $p$-grupo de orden $p^k$ con $k < n$ es nilpotente.
    
    \itshape Paso inductivo: \normalfont
    Sea $G$ un grupo de orden $p^n$. Sabemos, por la ecuación de clase para $p$-grupos, que el centro es no trivial: $Z(G) \neq \{e\}$.
    Consideremos el grupo cociente $G/Z(G)$.
    \begin{equation*}
        \left| \frac{G}{Z(G)} \right| = \frac{|G|}{|Z(G)|} = \frac{p^n}{p^m} = p^{n-m}.
    \end{equation*}
    Como $|Z(G)| \ge p$ (es decir $m \ge 1$), el orden del cociente es estrictamente menor que el de $G$.
    Por hipótesis de inducción, $G/Z(G)$ es nilpotente.
    
    Finalmente, aplicamos la observación demostrada previamente: si $G/Z(G)$ es nilpotente, entonces $G$ es nilpotente.
\end{proof}

\begin{lema}
    Sea $G$ un grupo finito.
    \begin{enumerate}
        \item Si $G$ es nilpotente y $H$ es un subgrupo propio de $G$ ($H \lneq G$), entonces $H$ es un subgrupo propio de su normalizador. Es decir:
        \begin{equation*}
            H \lneq G \implies H \lneq N_G(H).
        \end{equation*}
        \item Para cualquier subgrupo de Sylow $P$ de $G$, se cumple que $N_G(N_G(P)) = N_G(P)$.
    \end{enumerate}
\end{lema}

\begin{proof}
    \itshape Demostración de 2: \normalfont
    Sea $P \in \text{Syl}_p(G)$. Siempre se tiene que $N_G(P) \le N_G(N_G(P))$. Veamos la otra contención.
    Sea $x \in N_G(N_G(P))$. Por definición de normalizador:
    \begin{equation*}
        x N_G(P) x^{-1} = N_G(P).
    \end{equation*}
    Como $P \le N_G(P)$, al conjugar por $x$ tenemos que $xPx^{-1} \le x N_G(P) x^{-1} = N_G(P)$.
    
    Notemos que $P$ es un $p$-subgrupo de Sylow de $N_G(P)$. Más aún, como $P \triangleleft N_G(P)$, $P$ es el \emph{único} $p$-subgrupo de Sylow de $N_G(P)$.
    
    Por otro lado, como $xPx^{-1}$ es un conjugado de un $p$-Sylow, también es un $p$-Sylow.
    Dado que ambos ($P$ y $xPx^{-1}$) son subgrupos de Sylow dentro del grupo $N_G(P)$, y solo hay uno, deben ser iguales:
    \begin{equation*}
        xPx^{-1} = P.
    \end{equation*}
    Esto implica por definición que $x \in N_G(P)$.
    Por lo tanto, $N_G(N_G(P)) \le N_G(P)$, y se da la igualdad.
    
    \itshape Demostración de 1: \normalfont
    Sea $G$ nilpotente y $H \lneq G$. Procedemos por inducción sobre el orden de $G$.
    
    Si $G$ es abeliano, entonces $N_G(H) = G$. Como $H \lneq G$, entonces $H \lneq N_G(H)$ trivialmente.
    
    Supongamos que $G$ no es abeliano pero es nilpotente. Sabemos que $Z(G) \neq \{e\}$.
    Distinguimos dos casos:
    
    \itshape Caso A: $Z(G) \not\le H$. \normalfont
    Sea $z \in Z(G)$ tal que $z \notin H$.
    Como $z$ está en el centro, conmuta con todo elemento de $G$, y en particular con todo $h \in H$.
    Por tanto, $zHz^{-1} = H$, lo que implica $z \in N_G(H)$.
    Como $z \notin H$ y $z \in N_G(H)$, el normalizador contiene elementos fuera de $H$.
    Así, $H \lneq N_G(H)$ (pues contiene a $\langle H, z \rangle$).
    
    \itshape Caso B: $Z(G) \le H$. \normalfont
    Consideremos el homomorfismo canónico $\pi: G \to G/Z(G)$.
    Sea $\bar{G} = G/Z(G)$ y $\bar{H} = H/Z(G)$.
    
    Como $G$ es nilpotente, $\bar{G}$ es nilpotente. Además, como $Z(G) \neq \{e\}$, tenemos $|\bar{G}| < |G|$.
    Como $H \lneq G$, entonces $\bar{H} \lneq \bar{G}$.
    Podemos aplicar la hipótesis de inducción al grupo cociente:
    \begin{equation*}
        \bar{H} \lneq N_{\bar{G}}(\bar{H}).
    \end{equation*}
    Esto significa que existe un elemento $\bar{g} = gZ(G) \in \bar{G}$ tal que $\bar{g} \notin \bar{H}$ pero $\bar{g} \in N_{\bar{G}}(\bar{H})$.
    La condición $\bar{g} \in N_{\bar{G}}(\bar{H})$ implica:
    \begin{equation*}
        (gZ(G)) \frac{H}{Z(G)} (gZ(G))^{-1} = \frac{H}{Z(G)}.
    \end{equation*}
    Operando en el cociente:
    \begin{equation*}
        \frac{gHg^{-1}}{Z(G)} = \frac{H}{Z(G)} \implies gHg^{-1} = H.
    \end{equation*}
    Esto implica que $g \in N_G(H)$.
    
    Por otro lado, como $\bar{g} \notin \bar{H}$, entonces $g \notin H$.
    Hemos encontrado un elemento $g \in N_G(H) \setminus H$.
    Por lo tanto, $H \lneq N_G(H)$.
\end{proof}

\ObservacionBox{}{
    En la observación anterior sobre la relación con el centro, también es válida la versión para grupos solubles.    
}

\begin{prop}
    Sea $G$ un grupo finito. Si $G/Z(G)$ es soluble, entonces $G$ es soluble.
\end{prop}

\begin{proof}
    Supongamos que $G/Z(G)$ es soluble. Entonces existe un $n \in \mathbb{N}$ tal que su serie derivada llega a la identidad:
    \begin{equation*}
        \left(\frac{G}{Z(G)}\right)^{(n)} = \{Z(G)\}.
    \end{equation*}
    Usando la propiedad de los cocientes para series derivadas ($(G/N)^{(n)} = \frac{G^{(n)}N}{N}$), tenemos:
    \begin{equation*}
        \frac{G^{(n)}Z(G)}{Z(G)} = \{Z(G)\}.
    \end{equation*}
    Esto implica que $G^{(n)} \subseteq Z(G)$.
    
    Ahora calculamos el siguiente término de la serie derivada de $G$:
    \begin{equation*}
        G^{(n+1)} = [G^{(n)}, G^{(n)}].
    \end{equation*}
    Como $G^{(n)} \subseteq Z(G)$, el conmutador de cualesquiera dos elementos de $G^{(n)}$ es el conmutador de dos elementos centrales. Los elementos del centro conmutan entre sí, por lo que:
    \begin{equation*}
        [x, y] = e \quad \forall x, y \in G^{(n)}.
    \end{equation*}
    Por lo tanto:
    \begin{equation*}
        G^{(n+1)} = \{e\}.
    \end{equation*}
    Concluimos que $G$ es soluble.
\end{proof}

\TeoremaBox{}{
    Sea $G$ un grupo finito. $G$ es nilpotente si y solo si $G$ es el producto directo de sus subgrupos de Sylow.
}

\begin{proof}
    \itshape Dirección $(\Leftarrow)$: \normalfont
    Supongamos que $G$ es el producto directo de sus subgrupos de Sylow, digamos $G = P_1 \times \dots \times P_k$.
    Sabemos que todo $p$-grupo finito es nilpotente (visto en la observación anterior).
    Por la proposición demostrada previamente, el producto directo finito de grupos nilpotentes es nilpotente.
    Por lo tanto, $G$ es nilpotente.

    \itshape Dirección $(\Rightarrow)$: \normalfont
    Supongamos que $G$ es nilpotente. Sea $P$ un $p$-subgrupo de Sylow de $G$.
    Queremos demostrar que $P \triangleleft G$.
    Recordemos el Lema probado anteriormente: para cualquier subgrupo de Sylow $P$, se cumple que $N_G(N_G(P)) = N_G(P)$.
    
    Supongamos por contradicción que $N_G(P) \neq G$, es decir, $N_G(P)$ es un subgrupo propio de $G$.
    Como $G$ es nilpotente, aplicamos la propiedad de los normalizadores para subgrupos propios ($H \lneq N_G(H)$) tomando $H = N_G(P)$.
    Entonces:
    \begin{equation*}
        N_G(P) \lneq N_G(N_G(P)).
    \end{equation*}
    Pero esto contradice el lema que afirma $N_G(N_G(P)) = N_G(P)$.
    La única posibilidad es que $N_G(P) = G$.
    Esto implica que $P \triangleleft G$.
    
    Como esto vale para todo subgrupo de Sylow de $G$ (para cada primo que divide el orden), todos los subgrupos de Sylow son normales.
    Sabemos que si todos los subgrupos de Sylow de un grupo finito son normales, entonces el grupo es el producto directo de ellos.
    Concluimos que $G$ es el producto directo de sus subgrupos de Sylow.
\end{proof}

\CorolarioBox{}{
    Si $G$ es nilpotente finito y $m$ divide a $|G|$, entonces existe un subgrupo $H \le G$ tal que $|H| = m$.
}

\begin{proof}
    Por el teorema anterior, $G \cong P_1 \times \dots \times P_k$, donde cada $P_i$ es un $p_i$-grupo.
    Sea $m = p_1^{\alpha_1} \dots p_k^{\alpha_k}$ la descomposición prima de $m$.
    Como cada $P_i$ es un $p$-grupo, satisface el recíproco del Teorema de Lagrange (para $p$-grupos siempre existen subgrupos de cualquier orden que divida al total).
    Entonces, para cada $i$, existe $H_i \le P_i$ tal que $|H_i| = p_i^{\alpha_i}$.
    Tomando $H = H_1 \times \dots \times H_k$, tenemos un subgrupo de $G$ con orden:
    \begin{equation*}
        |H| = |H_1| \cdots |H_k| = p_1^{\alpha_1} \dots p_k^{\alpha_k} = m.
    \end{equation*}
\end{proof}

\TeoremaBox{Caracterizaciones de Nilpotencia}{
    Sea $G$ un grupo finito. Las siguientes afirmaciones son equivalentes:
    \begin{enumerate}
        \item $G$ es nilpotente (es decir, $G_{(n)} = \{e\}$ para algún $n$).
        \item Para todo subgrupo propio $H \lneq G$, se cumple $H \lneq N_G(H)$.
        \item La serie central ascendente alcanza al grupo, es decir, $C_{(n)} = G$ para algún $n$.
    \end{enumerate}
}

\begin{proof}
    Ya hemos demostrado $(1) \implies (2)$ en el lema de la página 27 y usamos $(2) \implies (1)$ implícitamente en el teorema de Sylow (vía la descomposición en producto directo). Ahora nos enfocaremos en la equivalencia con la serie central ascendente $(1) \iff (3)$.

    \itshape Demostración de $(1) \implies (3)$: \normalfont
    Sea $G$ nilpotente. Sea $n$ el mínimo entero tal que $G_{(n)} = \{e\}$.
    Demostraremos por inducción sobre $j$ que $G_{(n-j)} \subseteq C_{(j)}$ para todo $0 \le j \le n$.
    
    Para $j=0$: $G_{(n)} = \{e\} = C_{(0)}$. La base es cierta.
    
    Supongamos que $G_{(n-j)} \subseteq C_{(j)}$. Queremos ver que $G_{(n-(j+1))} \subseteq C_{(j+1)}$.
    Recordemos la definición de la serie ascendente mediante el conmutador:
    \begin{equation*}
        C_{(j+1)} = \{ g \in G \mid [g, x] \in C_{(j)} \forall x \in G \}.
    \end{equation*}
    Tomemos un elemento arbitrario $y \in G_{(n-j-1)}$. Queremos ver que $y \in C_{(j+1)}$.
    Para verificar esto, tomamos cualquier $x \in G$ y calculamos el conmutador $[x, y]$.
    Por definición de la serie central descendente:
    \begin{equation*}
        [x, y] \in [G, G_{(n-j-1)}] = G_{(n-j)}.
    \end{equation*}
    Por hipótesis de inducción, sabemos que $G_{(n-j)} \subseteq C_{(j)}$.
    Por lo tanto, $[x, y] \in C_{(j)}$.
    Esto implica, por la definición de $C_{(j+1)}$, que $y \in C_{(j+1)}$.
    
    Así hemos probado que $G_{(n-j-1)} \subseteq C_{(j+1)}$.
    Tomando $j=n$, obtenemos:
    \begin{equation*}
        G_{(0)} \subseteq C_{(n)}.
    \end{equation*}
    Como $G_{(0)} = G$, tenemos $G \subseteq C_{(n)}$, lo que implica $C_{(n)} = G$.

    \itshape Demostración de $(3) \implies (1)$: \normalfont
    Supongamos que $C_{(n)} = G$ para algún $n$.
    Demostraremos por inducción que $G_{(j)} \subseteq C_{(n-j)}$.
    
    Para $j=0$: $G_{(0)} = G = C_{(n)}$. La base es cierta.
    
    Supongamos que $G_{(j)} \subseteq C_{(n-j)}$. Queremos ver qué sucede con $G_{(j+1)}$.
    \begin{equation*}
        G_{(j+1)} = [G, G_{(j)}].
    \end{equation*}
    Por hipótesis de inducción, $G_{(j)} \subseteq C_{(n-j)}$. Usando la monotonía del conmutador:
    \begin{equation*}
        G_{(j+1)} \subseteq [G, C_{(n-j)}].
    \end{equation*}
    Recordemos la propiedad del término de la serie ascendente: $C_{(n-j)}/C_{(n-j-1)} = Z(G/C_{(n-j-1)})$.
    Esto significa que para todo $c \in C_{(n-j)}$ y $g \in G$, el conmutador $[g, c]$ pertenece a $C_{(n-j-1)}$.
    Por lo tanto:
    \begin{equation*}
        [G, C_{(n-j)}] \subseteq C_{(n-j-1)}.
    \end{equation*}
    Combinando las inclusiones:
    \begin{equation*}
        G_{(j+1)} \subseteq C_{(n-(j+1))}.
    \end{equation*}
    Tomando $j=n$, obtenemos:
    \begin{equation*}
        G_{(n)} \subseteq C_{(0)} = \{e\}.
    \end{equation*}
    Por lo tanto, $G_{(n)} = \{e\}$, y $G$ es nilpotente.
\end{proof}

 En la demostración anterior se verificó la relación inversa entre las series. Específicamente:
    \begin{itemize}
        \item Si la serie descendente termina en la identidad ($G_{(n)}=\{e\}$), entonces la ascendente cubre todo el grupo ($C_{(n)}=G$).
        \item Si la serie ascendente cubre todo el grupo, entonces la descendente termina en la identidad.
    \end{itemize}
    Esto unifica las definiciones de nilpotencia usando cualquiera de las dos series centrales.
